%========================================================================================
%  INTRODUZIONE  (si scrive per ULTIMA, quando i risultati sono definitivi)
%  1-2 pagine: contesto (PTA, evidenza del GWB 2023) -> problema (rumore delle pulsar,
%  forma spettrale uguale a quella del GWB -> degenerazioni) -> cosa ho fatto
%  (validazione Discovery vs Enterprise su J1640+2224, simulazioni, iniezioni di GWB,
%  scattering) -> risultati principali in 2-3 frasi -> struttura della tesi capitolo per capitolo.
%========================================================================================

\addcontentsline{toc}{chapter}{Introduzione}
\chapter*{Introduzione}
\markboth{Introduzione}{Introduzione}
\label{Introduzione}
\sloppy

\TODO{da scrivere alla fine}
